% sec10_1.tex — Section 10.1: Cointegrated Systems

Figure~10.1(a) plots the logs of quarterly real per capita aggregate personal dis\-posable
income and personal consumption expenditure for the United States from
1947:Q1 to 1998:Q1. Both series have linear trends and --- as will be verified in an
example below --- stochastic trends. However, these two I(1) series tend to move
together, suggesting that the difference is stationary. This is an example of \textbf{coin\-tegration}.
Cointegration relations abound in economics. In fact, many of the
variables we have examined in the book are cointegrated: prices of the same com\-modity
in two different countries (the difference should be stationary under the
weaker version of Purchasing Power Parity), long- and short-term interest rates
(even though they have trends, yield differential may be stationary), forward and
spot exchange rates, and so forth. Cointegration may characterize more than two
variables. For example, the existence of a stable money demand function implies
that a linear combination of the log of real money stock, the nominal interest rate,
and log aggregate income may be stationary even though each of the three variables
is I(1).

We start out this section by extending the definitions of univariate I(0) and I(1)
processes of the previous chapter to vector processes. The concept of cointegration
for multivariate I(1) processes will then be introduced formally.

\begin{figure}[htbp]
  \centering
  % Placeholder for Figure 10.1(a): Log Income and Consumption
  \includegraphics[width=0.9\textwidth,trim=50 350 50 350,clip]{figures/fig_10_1a.png}
  \caption{Log Income and Consumption}
  \label{fig:10.1a}
\end{figure}

\subsection*{Linear Vector I(0) and I(1) Processes}

Our definition of vector I(0) and I(1) processes follows Hamilton (1994) and
Johansen (1995). To introduce the multivariate extension of a zero-mean linear
univariate I(0) process, consider a zero-mean $n$-dimensional VMA (vector moving
average) process:
\begin{equation}
  \underset{(n \times 1)}{\mathbf{u}_t}
  = \underset{(n \times n)}{\boldsymbol{\Psi}(L)} \;
    \underset{(n \times 1)}{\boldsymbol{\varepsilon}_t}, \quad
  \boldsymbol{\Psi}(L) \equiv \mathbf{I}_n + \boldsymbol{\Psi}_1 L
    + \boldsymbol{\Psi}_2 L^2 + \cdots,
  \label{eq:10.1.1} \tag{10.1.1}
\end{equation}
where $\boldsymbol{\varepsilon}_t$ is i.i.d.\ with
\begin{equation}
  \E(\boldsymbol{\varepsilon}_t) = \mathbf{0}, \quad
  \E(\boldsymbol{\varepsilon}_t \boldsymbol{\varepsilon}_t')
  = \underset{(n \times n)}{\boldsymbol{\Omega}}, \quad
  \boldsymbol{\Omega} \text{ positive definite}.
  \label{eq:10.1.2} \tag{10.1.2}
\end{equation}
(The error process $\boldsymbol{\varepsilon}_t$ could be a vector Martingale Difference Sequence satisfying
certain properties, but we will not allow for this degree of generality.) As in the
univariate case, we impose two restrictions. The first is one-summability:\footnote{%
A sequence of matrices $\{\boldsymbol{\Psi}_j\}$ is said to be one-summable if
$\{|\boldsymbol{\Psi}_{j,k\ell}|\}$ is one-summable (i.e.,
$\sum_{j=0}^{\infty} j\,|\boldsymbol{\Psi}_{j,k\ell}| < \infty$)
for all $k, \ell = 1, 2, \ldots, n$, where $\boldsymbol{\Psi}_{j,k\ell}$ is the
$(k, \ell)$ element of the $n \times n$ matrix $\boldsymbol{\Psi}_j$.}
\begin{equation}
  \{\boldsymbol{\Psi}_j\} \text{ is one-summable}.
  \label{eq:10.1.3} \tag{10.1.3}
\end{equation}
Since $\{\boldsymbol{\Psi}_j\}$ is absolutely summable when one-summability holds, the multivariate
version of Proposition~6.1(d) means that $\mathbf{u}_t$ is (strictly) stationary and ergodic. The
second restriction on the zero-mean VMA process is
\begin{equation}
  \boldsymbol{\Psi}(1) \neq \underset{(n \times n)}{\mathbf{0}}
  \quad (n \times n \text{ matrix of zeros}).
  \label{eq:10.1.4} \tag{10.1.4}
\end{equation}
That is, at least one element of the $n \times n$ matrix $\boldsymbol{\Psi}(1)$ is not zero. This is a natu\-ral
multivariate extension of condition (9.2.3b) in the definition of univariate I(0)
processes on page~564.

A (linear) zero-mean $n$-dimensional \textbf{vector I(0) process} or a (linear) \textbf{zero-mean
I(0) system} is a VMA process satisfying (10.1.1)--(10.1.4). Given a zero-mean
I(0) system $\mathbf{u}_t$, a (linear) \textbf{vector I(0) process} or a (linear) \textbf{I(0) system} can
be written as
\[
  \boldsymbol{\delta} + \mathbf{u}_t,
\]
where $\boldsymbol{\delta}$ is an $n$-dimensional vector representing the mean of the stationary process.
Using the formula (6.5.6) with (6.3.15), the long-run covariance matrix of the I(0)
system is
\begin{equation}
  \boldsymbol{\Psi}(1) \boldsymbol{\Omega} \boldsymbol{\Psi}(1)'.
  \label{eq:10.1.5} \tag{10.1.5}
\end{equation}
Since $\boldsymbol{\Omega}$ is positive definite and $\boldsymbol{\Psi}(1)$ satisfies (10.1.4), at least one of the diagonal
elements of this long-run variance matrix is positive, implying that at least one of
the elements of an I(0) system is individually I(0) (i.e., I(0) as a univariate process).
We do not necessarily require $\boldsymbol{\Psi}(1)$ to be nonsingular. In fact, accommodating its
singularity is the whole point of the theory of cointegration. Consequently, the
long-run variance matrix, too, can be singular.

\begin{defbox}
\textbf{Example 10.1\quad (A bivariate I(0) system):}\quad Consider the following bivariate
first-order VMA process:
\begin{equation}
  \begin{cases}
    u_{1t} = \varepsilon_{1t} - \varepsilon_{1,t-1} + \gamma \varepsilon_{2,t-1}, \\
    u_{2t} = \varepsilon_{2t},
  \end{cases}
  \quad \text{or} \quad
  \mathbf{u}_t = \boldsymbol{\varepsilon}_t + \boldsymbol{\Psi}_1 \boldsymbol{\varepsilon}_{t-1},
  \label{eq:10.1.6} \tag{10.1.6}
\end{equation}
where
\begin{equation}
  \boldsymbol{\varepsilon}_t
  = \begin{bmatrix} \varepsilon_{1t} \\ \varepsilon_{2t} \end{bmatrix}, \quad
  \boldsymbol{\Psi}_1
  = \begin{bmatrix} -1 & \gamma \\ 0 & 0 \end{bmatrix}.
  \label{eq:10.1.7} \tag{10.1.7}
\end{equation}
Since this is a finite-order VMA, the one-summability condition is trivially
satisfied. Requirement (10.1.4) is satisfied because
\begin{equation}
  \boldsymbol{\Psi}(1) = \mathbf{I}_2 + \boldsymbol{\Psi}_1
  = \begin{bmatrix} 0 & \gamma \\ 0 & 1 \end{bmatrix}
  \neq \underset{(2 \times 2)}{\mathbf{0}}.
  \label{eq:10.1.8} \tag{10.1.8}
\end{equation}
So this example fits our definition of I(0) systems. If $\gamma = 0$, then the first
element, $u_{1t}$, is actually I($-1$) as a univariate process.
\end{defbox}

\medskip

The $n$-\textbf{dimensional I(1) system} $\mathbf{y}_t$ associated with a zero-mean I(0) system $\mathbf{u}_t$
can be defined by the relation
\begin{equation}
  \Delta \mathbf{y}_t = \boldsymbol{\delta} + \mathbf{u}_t.
  \label{eq:10.1.9} \tag{10.1.9}
\end{equation}
So the mean of $\Delta \mathbf{y}_t$ is $\boldsymbol{\delta}$. Since not every element of the associated zero-mean
I(0) process $\mathbf{u}_t$ is required to be individually I(0), some of the elements of the I(1)
system $\mathbf{y}_t$ may not be individually I(1).\footnote{%
In many textbooks and also in Engle and Granger (1987), all elements of an I(1) system are assumed to be
individually I(1), but that assumption is not needed for the development of cointegration theory. This point is
emphasized in, e.g., L\"{u}tkepohl (1993, pp.~352--353) and Johansen (1995, Chapter~3).}
Substituting (10.1.1) into this, we obtain
\begin{equation}
  \Delta \mathbf{y}_t = \boldsymbol{\delta} + \boldsymbol{\Psi}(L) \boldsymbol{\varepsilon}_t.
  \label{eq:10.1.10} \tag{10.1.10}
\end{equation}
This is called the \textbf{vector moving average (VMA) representation} of an I(1)
system. In levels, $\mathbf{y}_t$ can be written as
\begin{equation}
  \mathbf{y}_t = \mathbf{y}_0 + \boldsymbol{\delta} \cdot t
    + \mathbf{u}_1 + \mathbf{u}_2 + \cdots + \mathbf{u}_t,
  \label{eq:10.1.11} \tag{10.1.11}
\end{equation}
or, written out in full,
\[
  \begin{bmatrix} y_{1t} \\ y_{2t} \\ \vdots \\ y_{nt} \end{bmatrix}
  = \begin{bmatrix} y_{1,0} \\ y_{2,0} \\ \vdots \\ y_{n,0} \end{bmatrix}
  + \begin{bmatrix} \delta_1 \cdot t \\ \delta_2 \cdot t \\ \vdots \\ \delta_n \cdot t \end{bmatrix}
  + \begin{bmatrix} u_{1,1} \\ u_{2,1} \\ \vdots \\ u_{n,1} \end{bmatrix}
  + \begin{bmatrix} u_{1,2} \\ u_{2,2} \\ \vdots \\ u_{n,2} \end{bmatrix}
  + \cdots
  + \begin{bmatrix} u_{1t} \\ u_{2t} \\ \vdots \\ u_{nt} \end{bmatrix}.
\]
Regarding the initial value $\mathbf{y}_0$, we assume either that it is a vector of fixed constants
or that it is random but independent of $\boldsymbol{\varepsilon}_t$ for all~$t$.

\begin{defbox}
\textbf{Example 10.2\quad (A bivariate I(1) system):}\quad An I(1) system whose associated
zero-mean I(0) process is the bivariate process in Example~10.1 is
\begin{equation}
  \begin{cases}
    \Delta y_{1t} = \delta_1 + \varepsilon_{1t} - \varepsilon_{1,t-1}
      + \gamma \varepsilon_{2,t-1}, \\
    \Delta y_{2t} = \delta_2 + \varepsilon_{2t},
  \end{cases}
  \quad \text{or} \quad
  \Delta \mathbf{y}_t = \boldsymbol{\delta} + \boldsymbol{\varepsilon}_t
    + \boldsymbol{\Psi}_1 \boldsymbol{\varepsilon}_{t-1},
  \label{eq:10.1.12} \tag{10.1.12}
\end{equation}
where $\boldsymbol{\Psi}_1$ is given in (10.1.7). In levels,
\begin{equation}
  \begin{cases}
    y_{1t} = y_{1,0} + \delta_1 \cdot t + (\varepsilon_{1t} - \varepsilon_{1,0})
      + \gamma \,(\varepsilon_{2,0} + \varepsilon_{2,1} + \cdots + \varepsilon_{2,t-1}), \\
    y_{2t} = y_{2,0} + \delta_2 \cdot t + (\varepsilon_{2,1} + \varepsilon_{2,2}
      + \cdots + \varepsilon_{2t}).
  \end{cases}
  \label{eq:10.1.13} \tag{10.1.13}
\end{equation}
The second element of $\mathbf{y}_t$ is I(1) as a univariate process. If $\gamma = 0$, then
the first element $y_{1t}$ is trend stationary because the deviation from the trend
function $(y_{1,0} - \varepsilon_{1,0}) + \delta_1 \cdot t$ is stationary (actually i.i.d.). Otherwise $y_{1t}$ is I(1).
\end{defbox}

\subsection*{The Beveridge-Nelson Decomposition}

For an I(1) system whose associated zero-mean I(0) process is $\mathbf{u}_t$ satisfying
(10.1.1)--(10.1.4), it is easy to obtain the Beveridge-Nelson (BN) decomposition.
Recall from Section~9.2 that the univariate BN decomposition is based on the
rewriting of the MA filter as (9.2.5) on page~564. The obvious multivariate version
of it is
\begin{align}
  \underset{(n \times n)}{\boldsymbol{\Psi}(L)}
  &= \underset{(n \times n)}{\boldsymbol{\Psi}(1)}
    + (1 - L)\underset{(n \times n)}{\boldsymbol{\alpha}(L)}, \quad
  \boldsymbol{\alpha}(L) = \sum_{j=0}^{\infty} \boldsymbol{\alpha}_j L^j,
  \notag \\
  \underset{(n \times n)}{\boldsymbol{\alpha}_j}
  &= -(\boldsymbol{\Psi}_{j+1} + \boldsymbol{\Psi}_{j+2} + \cdots).
  \label{eq:10.1.14} \tag{10.1.14}
\end{align}
So the multivariate analogue of (9.2.6) on page~565 is
\begin{equation}
  \mathbf{u}_t = \boldsymbol{\Psi}(1) \boldsymbol{\varepsilon}_t
    + \boldsymbol{\eta}_t - \boldsymbol{\eta}_{t-1}, \quad
  \underset{(n \times 1)}{\boldsymbol{\eta}_t}
    \equiv \boldsymbol{\alpha}(L) \boldsymbol{\varepsilon}_t.
  \label{eq:10.1.15} \tag{10.1.15}
\end{equation}
Since $\boldsymbol{\Psi}(L)$ is one-summable, $\boldsymbol{\alpha}(L)$ is absolutely summable and $\boldsymbol{\eta}_t$ is a well-defined
covariance-stationary process. Substituting this into (10.1.11), we obtain
the multivariate version of the BN decomposition:
\begin{equation}
  \mathbf{y}_t = \boldsymbol{\delta} \cdot t
    + \boldsymbol{\Psi}(1)(\boldsymbol{\varepsilon}_1 + \boldsymbol{\varepsilon}_2
      + \cdots + \boldsymbol{\varepsilon}_t)
    + \boldsymbol{\eta}_t
    + (\mathbf{y}_0 - \boldsymbol{\eta}_0).
  \label{eq:10.1.16} \tag{10.1.16}
\end{equation}
As in the univariate case, the I(1) system $\mathbf{y}_t$ is decomposed into a linear trend
$\boldsymbol{\delta} \cdot t$,
a stochastic trend $\boldsymbol{\Psi}(1)(\boldsymbol{\varepsilon}_1 + \boldsymbol{\varepsilon}_2 + \cdots + \boldsymbol{\varepsilon}_t)$,
a stationary component $\boldsymbol{\eta}_t$, and the
initial condition $\mathbf{y}_0 - \boldsymbol{\eta}_0$. By construction, $\boldsymbol{\eta}_0$ is a random variable. So unless
$\mathbf{y}_0$ is perfectly correlated with $\boldsymbol{\eta}_0$, the initial condition is random.

\begin{defbox}
\textbf{Example 10.2 (continued):}\quad For the bivariate I(1) system in Example~10.2,
the matrix version of $\boldsymbol{\alpha}(L)$ equals $-\boldsymbol{\Psi}_1$ so that the stationary component $\boldsymbol{\eta}_t$
in the BN decomposition is
\begin{equation}
  \boldsymbol{\eta}_t
  = \begin{bmatrix} \varepsilon_{1t} - \gamma \varepsilon_{2t} \\ 0 \end{bmatrix}.
  \label{eq:10.1.17} \tag{10.1.17}
\end{equation}
(It should be easy for you to verify this from the matrix version of (9.2.5).)
The two-dimensional stochastic trend is written as
\begin{equation}
  \boldsymbol{\Psi}(1)(\boldsymbol{\varepsilon}_1 + \boldsymbol{\varepsilon}_2
    + \cdots + \boldsymbol{\varepsilon}_t)
  = \begin{bmatrix} 0 & \gamma \\ 0 & 1 \end{bmatrix}
    \begin{bmatrix} \sum_{s=1}^{t} \varepsilon_{1s} \\
                     \sum_{s=1}^{t} \varepsilon_{2s} \end{bmatrix}
  = \begin{bmatrix} \gamma \cdot \sum_{s=1}^{t} \varepsilon_{2s} \\
                     \sum_{s=1}^{t} \varepsilon_{2s} \end{bmatrix}.
  \label{eq:10.1.18} \tag{10.1.18}
\end{equation}
Thus, the two-dimensional stochastic trend is driven by one common stochas\-tic
trend, $\sum_{s=1}^{t} \varepsilon_{2s}$. This is an example of the ``common trend'' representation
of Stock and Watson (1993) (see Review Question~2 below).
\end{defbox}

\subsection*{Cointegration Defined}

The I(1) system $\mathbf{y}_t$ is not stationary because it contains a stochastic trend
$\boldsymbol{\Psi}(1) \cdot (\boldsymbol{\varepsilon}_1 + \cdots + \boldsymbol{\varepsilon}_t)$,
and, if $\boldsymbol{\delta} \neq \mathbf{0}$, a deterministic trend $\boldsymbol{\delta} \cdot t$.
The deterministic and stochastic trends, however, may disappear if we take a suitable linear combina\-tion
of the elements of $\mathbf{y}_t$. To pursue this idea,\footnote{%
The idea was originally suggested by Granger (1981). See also Aoki (1968) and Box and Tiao (1977).}
premultiply both sides of the BN
decomposition (10.1.16) by some conformable vector $\mathbf{a}$ to obtain
\begin{equation}
  \mathbf{a}' \mathbf{y}_t = \mathbf{a}' \boldsymbol{\delta} \cdot t
    + \mathbf{a}' \boldsymbol{\Psi}(1)
      (\boldsymbol{\varepsilon}_1 + \boldsymbol{\varepsilon}_2 + \cdots + \boldsymbol{\varepsilon}_t)
    + \mathbf{a}' \boldsymbol{\eta}_t
    + \mathbf{a}'(\mathbf{y}_0 - \boldsymbol{\eta}_0).
  \label{eq:10.1.19} \tag{10.1.19}
\end{equation}
If $\mathbf{a}$ satisfies
\begin{equation}
  \mathbf{a}' \boldsymbol{\Psi}(1) = \underset{(1 \times n)}{\mathbf{0}'},
  \label{eq:10.1.20} \tag{10.1.20}
\end{equation}
then the stochastic trend is eliminated and $\mathbf{a}' \mathbf{y}_t$ becomes
\begin{equation}
  \mathbf{a}' \mathbf{y}_t
  = \mathbf{a}' \boldsymbol{\delta} \cdot t
    + \mathbf{a}' \boldsymbol{\eta}_t
    + \mathbf{a}'(\mathbf{y}_0 - \boldsymbol{\eta}_0).
  \label{eq:10.1.21} \tag{10.1.21}
\end{equation}
Strictly speaking, this process is not trend stationary because the initial condition
$\mathbf{a}'(\mathbf{y}_0 - \boldsymbol{\eta}_0)$ can be correlated with subsequent values of
$\mathbf{a}' \boldsymbol{\eta}_t$ (see Review Ques\-tion~3 below for an illustration). The process would be trend stationary if, for
example, the initial value $\mathbf{y}_0$ were such that
$\mathbf{a}'(\mathbf{y}_0 - \boldsymbol{\eta}_0) = 0$.

We are now ready to define cointegration.\footnote{%
Our definition is the same as Definition 3.4 in Johansen (1995).}

\begin{definition}[Cointegration]
\label{def:10.1}
Let $\mathbf{y}_t$ be an $n$-dimensional I(1) system whose
associated zero-mean I(0) system $\mathbf{u}_t$ satisfies (10.1.1)--(10.1.4). We say that $\mathbf{y}_t$ is
\textbf{cointegrated} with an $n$-dimensional \textbf{cointegrating vector} $\mathbf{a}$ if $\mathbf{a} \neq \mathbf{0}$ and
$\mathbf{a}' \mathbf{y}_t$ can
be made trend stationary by a suitable choice of its initial value $\mathbf{y}_0$.
\end{definition}

Defining cointegration in this way, although dictated by logical consistency, does
\emph{not} necessarily mean that the theory of cointegration requires that the initial con\-dition
$\mathbf{y}_0$ be chosen as indicated in the definition. The process in (10.1.21) is not
stationary because $\boldsymbol{\eta}_t$ and $\boldsymbol{\eta}_0$ are correlated. However, since
$\boldsymbol{\eta}_t = \boldsymbol{\alpha}(L) \boldsymbol{\varepsilon}_t$ and $\boldsymbol{\alpha}(L)$
is absolutely summable, $\boldsymbol{\eta}_t$ and $\boldsymbol{\eta}_0$ will become asymptotically independent as $t$
increases. In this sense, the process in (10.1.21) is ``asymptotically stationary,''
and asymptotic stationarity is all that is needed for estimation and inference for
cointegrated I(1) systems.\footnote{%
The distinction between stationarity and asymptotic stationarity is discussed in L\"{u}tkepohl (1993, pp.~348--349).}

\medskip

With cointegration thus defined, we can define the following related concepts.

\begin{itemize}
\item (Cointegration rank)\quad The \textbf{cointegrating rank} is the number of linearly
  independent cointegrating vectors, and the \textbf{cointegrating space} is the space
  spanned by the cointegrating vectors. As the preceding discussion shows, an
  $n$-dimensional vector $\mathbf{a}$ is a cointegrating vector if and only if (10.1.20) holds.
  Since there are $h$ linearly independent such vectors if the cointegration rank is
  $h$, it follows that
  \begin{equation}
    \rank\bigl[\underset{(n \times n)}{\boldsymbol{\Psi}(1)}\bigr] = n - h.
    \label{eq:10.1.22} \tag{10.1.22}
  \end{equation}
  Put differently, the cointegration rank $h$ equals $n - \rank[\boldsymbol{\Psi}(1)]$.

\item (Cointegration among subsets of $\mathbf{y}_t$)\quad At least one of the elements of a coin\-tegrating
  vector is not zero. Suppose, without loss of generality, that the first
  element of the cointegrating vector is not zero. Then we say that $y_{1t}$ (the first
  element of $\mathbf{y}_t$) is \textbf{cointegrated with} $\mathbf{y}_{2t}$ (the remaining $n - 1$ elements of $\mathbf{y}_t$) or
  that $y_{1t}$ is \textbf{part of a cointegrating relationship}. We can also define cointegra\-tion
  for a subset of $\mathbf{y}_t$. For example, the $n - 1$ variables in $\mathbf{y}_{2t}$ are not cointegrated
  if there does not exist an $(n - 1)$-dimensional vector $\mathbf{b} \neq \mathbf{0}$ such that (10.1.20)
  holds with $\mathbf{a}' = (0, \mathbf{b}')$. Therefore, $\mathbf{y}_{2t}$ is not cointegrated if and only if the last
  $n - 1$ rows of $\boldsymbol{\Psi}(1)$ are linearly independent.

\item (Stochastic cointegration)\quad Note that the deterministic trend $\mathbf{a}' \boldsymbol{\delta} \cdot t$ is not elimi\-nated
  from $\mathbf{a}' \mathbf{y}_t$ unless the cointegrating vector also satisfies
  \begin{equation}
    \mathbf{a}' \boldsymbol{\delta} = 0.
    \label{eq:10.1.23} \tag{10.1.23}
  \end{equation}
\end{itemize}

In most applications, a cointegrating vector eliminates both the stochastic and
deterministic trends in (10.1.19). That is, a vector $\mathbf{a}$ that satisfies (10.1.20) neces\-sarily
satisfies (10.1.23), so that $\mathbf{a}' \mathbf{y}_t$, which can now be written as
\begin{equation}
  \mathbf{a}' \mathbf{y}_t
  = \mathbf{a}' \boldsymbol{\eta}_t
    + \mathbf{a}'(\mathbf{y}_0 - \boldsymbol{\eta}_0),
  \label{eq:10.1.24} \tag{10.1.24}
\end{equation}
is stationary (not just trend stationary) for a suitable choice of $\mathbf{y}_0$. This implies
that $\boldsymbol{\delta}$ is a linear combination of the columns of $\boldsymbol{\Psi}(1)$, so
\begin{equation}
  \rank \begin{bmatrix}
    \boldsymbol{\delta} & \vdots & \boldsymbol{\Psi}(1) \\
    \multicolumn{3}{c}{\underset{(n \times (n+1))}{}}
  \end{bmatrix} = n - h.
  \label{eq:10.1.25} \tag{10.1.25}
\end{equation}
Unless otherwise noted, we will assume that (10.1.25) as well as (10.1.22) are
satisfied when the cointegration rank is~$h$. If we wish to describe the case where
a cointegrating vector eliminates the stochastic trend but not necessarily the de\-terministic
trend, we will use the term \textbf{stochastic cointegration}. Of course, if $\mathbf{y}_t$
does not contain a deterministic trend (i.e., if $\boldsymbol{\delta} = \mathbf{0}$), then there is no difference
between cointegration and stochastic cointegration.

\begin{defbox}
\textbf{Example 10.2 (continued):}\quad In the bivariate I(1) system (10.1.12) in Exam\-ple~10.2,
the matrix $\boldsymbol{\Psi}(1)$ is given in (10.1.8), so
\[
  \begin{bmatrix}
    \boldsymbol{\delta} & \vdots & \boldsymbol{\Psi}(1)
  \end{bmatrix}
  = \begin{bmatrix}
    \delta_1 & 0 & \gamma \\
    \delta_2 & 0 & 1
  \end{bmatrix}.
\]
The rank of $\boldsymbol{\Psi}(1)$ is 1, so the cointegration rank is $1$ ($= 2 - 1$). All cointe\-grating
vectors can be written as $(c, -c\gamma)'$, $c \neq 0$. The assumption that the
cointegrating vector also eliminates the deterministic trend can be written as
$c\delta_1 - c\gamma\delta_2 = 0$ or $\delta_1 = \gamma\delta_2$, which implies that the rank of
$[\boldsymbol{\delta} \vdots \boldsymbol{\Psi}(1)]$ shown
above is one.
\end{defbox}

\medskip

Quite a few implications for I(1) systems follow from the definition of cointe\-gration.
For example,

\begin{itemize}
\item ($h < n$)\quad For an $n$-dimensional I(1) system, the cointegration rank cannot be
  as large as~$n$. If it were $n$, then $\rank[\boldsymbol{\Psi}(1)] = 0$ and $\boldsymbol{\Psi}(1)$ would have to be a
  matrix of zeros, which is ruled out by requirement (10.1.4).

\item (Implications of the positive definiteness of the long-run variance of $\Delta \mathbf{y}_t$)\quad The
  long-run variance matrix of $\Delta \mathbf{y}_t$ is given by $\boldsymbol{\Psi}(1)\boldsymbol{\Omega}\boldsymbol{\Psi}(1)'$
  (see (10.1.5)), which
  is positive definite if and only if $\boldsymbol{\Psi}(1)$ is nonsingular. Therefore, $\mathbf{y}_t$ is not coin\-tegrated
  if and only if the long-run variance matrix of $\Delta \mathbf{y}_t$ is positive definite.\footnote{%
  This equivalence is specific to linear processes. In general, the positive definiteness of the long-run variance
  matrix of $\Delta \mathbf{y}_t$ is sufficient, but not always necessary, for $\mathbf{y}_t$ to be not cointegrated. See Phillips (1986, p.~321).}
  Since the long-run variance of each element of $\Delta \mathbf{y}_t$ is positive if the long-run
  variance matrix of $\Delta \mathbf{y}_t$ is positive definite, it follows that each element of $\mathbf{y}_t$ is
  individually I(1) (i.e., I(1) as a univariate process) if $\mathbf{y}_t$ is not cointegrated. The
  same is true for a subset of $\mathbf{y}_t$. For example, consider $\mathbf{y}_{2t}$, the last $n - 1$ elements
  of $\mathbf{y}_t$. The long-run variance matrix of $\Delta \mathbf{y}_{2t}$ is given by
  $\boldsymbol{\Psi}_2(1)\boldsymbol{\Omega}\boldsymbol{\Psi}_2(1)'$ where
  $\boldsymbol{\Psi}_2(1)$ is the last $n - 1$ rows of $\boldsymbol{\Psi}(1)$. It is positive definite if and only if the rows
  of $\boldsymbol{\Psi}_2(1)$ are linearly independent or $\mathbf{y}_{2t}$ is not cointegrated. In particular, each
  element of $\mathbf{y}_{2t}$ is individually I(1) if $\mathbf{y}_{2t}$ is not cointegrated.

\item Suppose that $y_{1t}$ is cointegrated with $\mathbf{y}_{2t}$. Then $\mathbf{y}_{2t}$ is not cointegrated if $h = 1$
  and cointegrated if $h > 1$. The reason is as follows. Cointegration of $y_{1t}$ with
  $\mathbf{y}_{2t}$ implies that there exists a cointegrating vector, call it $\mathbf{a}_1$, whose first element
  is not zero. If $h = 1$, then there should not exist an $(n - 1)$-dimensional vector
  $\mathbf{b} \neq \mathbf{0}$ such that $\mathbf{a}_2' = (0, \mathbf{b}')$ is a cointegrating vector, because $\mathbf{a}_1$ and $\mathbf{a}_2$ are
  linearly independent. A vector such as $\mathbf{a}_2$ can be found if $h > 1$.

\item (VAR in first differences)\quad If $\mathbf{y}_t$ is difference stationary (without drift), it is
  tempting to model it as a stationary finite-order VAR
  $\boldsymbol{\Phi}(L) \Delta \mathbf{y}_t = \boldsymbol{\varepsilon}_t$ where
  $\boldsymbol{\Phi}(L)$ is a matrix polynomial of degree $p$ such that all the roots of
  $|\boldsymbol{\Phi}(z)| = 0$
  are outside the unit circle. But then $\mathbf{y}_t$ cannot be cointegrated. The reason is
  as follows. If $\boldsymbol{\Phi}(L)$ satisfies the stationarity condition, the coefficient matrix
  sequence $\{\boldsymbol{\Psi}_j\}$ of its inverse
  $\boldsymbol{\Psi}(L) = \boldsymbol{\Phi}(L)^{-1}$ is bounded by a geometrically
  declining sequence (see Section~6.3). So $\boldsymbol{\Psi}(L)$ is one-summable and
  $\Delta \mathbf{y}_t = \boldsymbol{\Psi}(L) \boldsymbol{\varepsilon}_t$ is a VMA process satisfying (10.1.2) and (10.1.3). Furthermore,
  \[
    |\boldsymbol{\Psi}(1)| = |\boldsymbol{\Phi}(1)|^{-1} = |\boldsymbol{\Phi}(1)|^{-1} \neq 0.
  \]
  So $\boldsymbol{\Psi}(1)$ is nonsingular and the long-run variance matrix of $\Delta \mathbf{y}_t$ is positive
  definite.
\end{itemize}

\subsection*{Questions for Review}

\begin{enumerate}
\item Let $(y_{1t}, y_{2t})$ be as in Example~10.2 with $\gamma = 0$. Show that $\{y_{1t} + y_{2t}\}$ is I(1).
  \textbf{Hint:} You need to show that the long-run variance of $\{\Delta y_{1t} + \Delta y_{2t}\}$ is positive.

\item (Stock and Watson (1988) common-trend representation)\quad
  Let $\mathbf{w}_t \equiv \boldsymbol{\varepsilon}_1 + \cdots + \boldsymbol{\varepsilon}_t$ so that the BN decomposition is
  \[
    \mathbf{y}_t = \boldsymbol{\delta} \cdot t + \boldsymbol{\Psi}(1) \mathbf{w}_t
      + \boldsymbol{\eta}_t + (\mathbf{y}_0 - \boldsymbol{\eta}_0).
  \]
  A result from linear algebra states that

  if $\mathbf{C}$ is an $n \times n$ matrix of rank $n - h$, then there exists an $n \times n$ non\-singular
  matrix $\mathbf{G}$ and an $n \times (n - h)$ matrix $\mathbf{F}$ of full column rank
  such that
  \[
    \underset{(n \times n)}{\mathbf{C}} \, \underset{(n \times n)}{\mathbf{G}}
    = \begin{bmatrix}
        \underset{(n \times (n-h))}{\mathbf{F}} & \vdots
          & \underset{(n \times h)}{\mathbf{0}}
      \end{bmatrix}.
  \]

  Show that the permanent component $\boldsymbol{\Psi}(1) \mathbf{w}_t$ can be written as
  $\mathbf{F} \boldsymbol{\tau}_t$, where $\mathbf{F}$ is
  an $n \times (n - h)$ matrix of full column rank and $\boldsymbol{\tau}_t$ is an $(n - h)$-dimensional
  I(1) process (without drift). This is called the \textbf{common-trend representation},
  and $\boldsymbol{\tau}_t$ is the ($n - h$)-dimensional \textbf{common stochastic trend}.

\item Show that the I(1) system in Example~10.2 can be cointegrated but not trend
  stationary. \textbf{Hint:} Show that $\mathbf{a}'(\mathbf{y}_0 - \boldsymbol{\eta}_0) \neq 0$ for
  $\mathbf{a} = (1, -\gamma)'$ if
  $y_{1,0} - \varepsilon_{1,0} + \gamma \varepsilon_{2,0} \neq -\gamma y_{2,0}$.
\end{enumerate}
